\chapter{Images And Pixmaps}

Chapter 7 used images as painter targets. This chapter looks at images as data: loading files, probing metadata, reading and writing bytes, transforming raster content, moving pixels between \api{QImage} and \api{QPixmap}, and using the clipboard.

\api{QImage} is the preferred representation for direct pixel and byte-buffer work. \api{QPixmap} is best for widget-facing images, grabs, and display-oriented raster data.

\begin{longtable}{@{}p{0.24\textwidth}p{0.70\textwidth}@{}}
\toprule
Type & Use it for \\
\midrule
\api{QImage} & Pixel access, raw bytes, format conversion, offscreen drawing, file encoding, clipboard images. \\
\api{QPixmap} & Widget grabs, display-oriented images, pixmap brushes, conversion to and from \api{QImage}. \\
\api{QImageReader} & Probing and decoding image files or image data from an \api{IODevice}. \\
\api{QBuffer} & Treating an in-memory byte array as a readable or writable Qt device. \\
\api{QByteArray} & Passing copied binary payloads between Qt APIs and Crystal code. \\
\api{MimeData} & Clipboard and drag/drop payloads that combine text, HTML, images, and custom bytes. \\
\bottomrule
\end{longtable}

\section{Image Or Pixmap?}

Use \api{QImage} when your code needs to inspect pixels, edit pixels, encode bytes, decode from bytes, or render offscreen. It is the natural choice for import pipelines, thumbnails, generated previews, image-processing steps, and tests that assert exact pixel colors.

Use \api{QPixmap} when the image is tied to widgets. \api{Widget#grab} returns a pixmap, and pixmaps can be used for display snapshots, textured brushes, splash screens, and conversions back to \api{QImage} when pixel inspection is needed.

\begin{lstlisting}[style=crystal]
image = Qt6::QImage.new(256, 256, Qt6::ImageFormat::ARGB32)
image.fill(Qt6::Color.new(20, 24, 28))
image.set_pixel_color(12, 12, Qt6::Color.new(255, 80, 80))

thumbnail = image.scaled(96, 96, Qt6::AspectRatioMode::Keep, Qt6::TransformationMode::Smooth)
pixmap = thumbnail.to_pixmap

round_trip = pixmap.to_image
round_trip.pixel_color(12, 12)
\end{lstlisting}

\begin{figure}[htbp]
  \centering
  \screenshotimage{docs/book/images/images-pipeline.png}{width=0.92\textwidth,height=0.44\textheight,keepaspectratio}%
  \caption{A typical image workflow: bytes become a \api{QImage}, transforms produce derived rasters, and \api{QPixmap} sits at the widget boundary.}
\end{figure}

\section{Creating And Editing Images}

\api{QImage.new} creates an image with dimensions and an optional format. The common default is \api{ImageFormat::ARGB32}, which stores a full alpha channel.

\begin{lstlisting}[style=crystal]
image = Qt6::QImage.new(320, 180)
image.fill(Qt6::Color.new(255, 255, 255))

image.set_pixel_color(40, 32, Qt6::Color.new(220, 72, 60))
image.set_pixel_color(41, 32, Qt6::Color.new(220, 72, 60))

color = image.pixel_color(40, 32)
status_bar.show_message("Red #{color.red}, alpha #{color.alpha}", 1200)
\end{lstlisting}

For rectangular edits, prefer copy and painter operations over per-pixel loops when possible.

\begin{lstlisting}[style=crystal]
tile = image.copy(Qt6::Rect.new(32, 24, 96, 96))
tile = tile.convert_to_format(Qt6::ImageFormat::ARGB32Premultiplied)
tile.invert_pixels
\end{lstlisting}

The binding exposes common metadata:

\begin{lstlisting}[style=crystal]
image.width
image.height
image.size
image.format
image.depth
image.bytes_per_line
image.size_in_bytes
image.has_alpha_channel?
image.null?
\end{lstlisting}

\section{Raw Bytes}

\api{QImage.from_raw_data} copies raw pixel bytes into a new image. You provide the width, height, bytes per line, and format.

\begin{lstlisting}[style=crystal]
pixels = Bytes[
  30_u8, 20_u8, 10_u8, 255_u8,
  60_u8, 50_u8, 40_u8, 255_u8,
  90_u8, 80_u8, 70_u8, 255_u8,
  120_u8, 110_u8, 100_u8, 255_u8,
]

image = Qt6::QImage.from_raw_data(
  pixels,
  width: 2,
  height: 2,
  bytes_per_line: 8,
  format: Qt6::ImageFormat::ARGB32
)
\end{lstlisting}

For \api{ARGB32}, the in-memory byte order for little-endian systems is commonly blue, green, red, alpha. The example above reads back as red \api{10}, green \api{20}, blue \api{30} for the first pixel.

\begin{lstlisting}[style=crystal]
first = image.pixel_color(0, 0)
first.red   # 10
first.green # 20
first.blue  # 30
\end{lstlisting}

\api{QImage#bytes} returns a copied byte slice. That makes it safe to inspect or pass to ordinary Crystal code without keeping a borrowed Qt pointer alive.

\section{Scaling And Transforms}

\api{QImage} and \api{QPixmap} both support fixed-size scaling, aspect-ratio-preserving scaling, and transform copies.

\begin{lstlisting}[style=crystal]
preview = image.scaled(
  480,
  320,
  Qt6::AspectRatioMode::Keep,
  Qt6::TransformationMode::Smooth
)

wide = image.scaled_to_width(640, Qt6::TransformationMode::Smooth)
tall = image.scaled_to_height(480)
\end{lstlisting}

Choose the aspect-ratio mode by intent:

\begin{itemize}
\item \api{Keep}: fit inside a box.
\item \api{KeepByExpanding}: fill a box, then crop if needed.
\item \api{Ignore}: stretch to the requested dimensions.
\end{itemize}

\begin{lstlisting}[style=crystal]
transform = Qt6::QTransform.new
transform.rotate(90)

rotated = image.transformed(transform, Qt6::TransformationMode::Smooth)
mirrored = image.mirrored(horizontal: true, vertical: false)
swapped = image.rgb_swapped
\end{lstlisting}

\section{Loading Files}

For direct loading, use \api{QImage.from_file}, \api{QPixmap.from_file}, or \api{load} on an existing object.

\begin{lstlisting}[style=crystal]
image = Qt6::QImage.from_file("assets/terrain.png")
unless image.null?
  preview.pixmap = image.scaled_to_width(320).to_pixmap
end

pixmap = Qt6::QPixmap.from_file("assets/icon.png")
\end{lstlisting}

For editor-style workflows, prefer \api{QImageReader}. It can probe a file before decoding it, report the image size, expose the detected format, enable auto transforms, and provide an error string when decoding fails.

\begin{lstlisting}[style=crystal]
reader = Qt6::QImageReader.new(path)
reader.auto_transform = true

if reader.can_read?
  size = reader.size
  status_bar.show_message("#{reader.format.upcase} #{size.width}x#{size.height}", 1500)
  image = reader.read
else
  status_bar.show_message(reader.error_string, 2200)
end
\end{lstlisting}

You can also decode into an existing image instance:

\begin{lstlisting}[style=crystal]
destination = Qt6::QImage.new(1, 1)
if reader.read(destination)
  preview.pixmap = destination.to_pixmap
end
\end{lstlisting}

\section{Encoding And In-Memory IO}

\api{save} writes to a file path. \api{save_to_data} returns encoded bytes for a named format such as \api{PNG}.

\begin{lstlisting}[style=crystal]
image.save("preview.png")

png_bytes = image.save_to_data("PNG")
File.write("preview-copy.png", png_bytes)
\end{lstlisting}

\api{QBuffer} lets Qt APIs read from or write to memory through the same \api{IODevice} interface they use for files.

\begin{lstlisting}[style=crystal]
buffer = Qt6::QBuffer.new
buffer.open(Qt6::IODeviceOpenMode::ReadWrite)

image.save(buffer, "PNG")
buffer.seek(0)

decoded = Qt6::QImage.new(1, 1)
decoded.load(buffer, "PNG")
buffer.close
\end{lstlisting}

That pattern is useful when you want to round-trip generated images without touching disk, store encoded images in a database, send images over the network, or feed encoded bytes into \api{QImageReader}.

\begin{lstlisting}[style=crystal]
buffer.seek(0)
reader = Qt6::QImageReader.new(buffer, "png")
image = reader.read if reader.can_read?
\end{lstlisting}

\begin{figure}[htbp]
  \centering
  \screenshotimage{docs/book/images/images-io-clipboard.png}{width=0.92\textwidth,height=0.44\textheight,keepaspectratio}%
  \caption{Images can move from files through \api{QImageReader}, into memory buffers, and onward through \api{MimeData} and the clipboard.}
\end{figure}

\section{QByteArray}

\api{QByteArray} is the small binary payload wrapper used by \api{QBuffer}, \api{IODevice}, and some MIME flows.

\begin{lstlisting}[style=crystal]
payload = Qt6::QByteArray.new(Bytes[1_u8, 7_u8, 9_u8])
payload.size
payload.bytes
payload.clear
\end{lstlisting}

Many APIs return copied bytes, either as \api{Bytes} or as a \api{QByteArray} that can be converted to \api{Bytes}. Treat those as ordinary Crystal data after the call returns.

\section{Widget Grabs}

Any widget can be captured with \api{grab}. The result is a \api{QPixmap}, which can be saved directly or converted to a \api{QImage}.

\begin{lstlisting}[style=crystal]
pixmap = preview.grab
pixmap.save("preview-widget.png")

image = pixmap.to_image
sample = image.pixel_color(10, 10)
\end{lstlisting}

This is handy for export commands, documentation screenshots, tests, and preview thumbnails. If you need the widget to repaint before grabbing it, call \api{update}, process events, and then grab.

\section{Clipboard And MIME Data}

The application clipboard supports text, images, pixmaps, and MIME payloads.

\begin{lstlisting}[style=crystal]
clipboard = Qt6.application.clipboard

clipboard.image = image
if clipboard.has_image?
  pasted = clipboard.image
  preview.pixmap = pasted.to_pixmap
end

clipboard.pixmap = preview.grab
\end{lstlisting}

Use \api{MimeData} when clipboard or drag/drop data needs to carry multiple representations at once.

\begin{lstlisting}[style=crystal]
mime = Qt6::MimeData.new
mime.text = "Layer preview"
mime.html = "<b>Layer preview</b>"
mime.image = image
mime.set_data("application/x-crystal-qt6-layer", Bytes[1_u8, 7_u8, 9_u8])

clipboard.mime_data = mime
\end{lstlisting}

On the receiving side, check for the representation you want before reading it.

\begin{lstlisting}[style=crystal]
if data = clipboard.mime_data
  if data.has_image?
    preview.pixmap = data.image.to_pixmap
  elsif data.has_text?
    status_bar.show_message(data.text, 1500)
  end
end
\end{lstlisting}

\section{Practical Guidelines}

Keep import code defensive. Check \api{null?}, \api{can_read?}, or return values from \api{load} and \api{save}; surface \api{QImageReader#error_string} when loading fails.

Use \api{QImage} as the working format for generated images, pixel tests, and encoded data. Convert to \api{QPixmap} at the UI boundary.

Use \api{TransformationMode::Smooth} for user-facing thumbnails and previews. Use \api{Fast} when speed matters more than filtering quality.

Prefer \api{QImageReader} before decoding large user-selected assets. It lets the UI show useful metadata and reject unsupported files gracefully.

Use \api{QBuffer} for in-memory round trips and \api{MimeData} for clipboard or drag/drop flows that need more than one representation.
